% =============================================================================
\section{From trajectories to couplings}\label{sec:inference}

\subsection{Fitting and choosing the depth}\label{sec:estimation}

For walkers $i=1,\dots,n$ with flights $x_{i,t}$, the field of Eq.~\eqref{eq:causal} becomes $\theta_{i,t}=h_i+\sum_\tau J_\tau x_{i,t-\tau}$, with one field per walker and shared couplings. The log-likelihood stays concave, and Newton's method finds its maximum in a few iterations (Appendix~\ref{app:estimation}). The cost is negligible. For one larval condition (6 to 59 larvae, 349 to 4240 flights) the whole scan of depths $K=0,\dots,6$, with one field per larva, takes 0.1--0.6\,s on a single core of a desktop computer; pooling all 241 larvae, with 241 fields and up to six couplings for 20\,948 flights, takes 5\,s. The same machinery accommodates fields that change every $B$ flights, which absorb slow drifts of the exponent and leave to the couplings only the short-range memory, and known covariates, such as the distance from a wall. Short tracks bias the couplings of a fit with one field per walker by an amount of order $1/T_i$~\cite{Nickell1981}, which we measure on surrogates in which the flights of each walker are permuted.

The depth is chosen by comparing models with $K=0,\dots,K_{\max}$, fitted on the same flights, through an information criterion, $\mathrm{AIC}=-2l+2k$ or $\mathrm{BIC}=-2l+k\ln n_{\rm o}$, with $k$ parameters and $n_{\rm o}$ predicted flights~\cite{Akaike1974,Schwarz1978}. The choice is not innocent. On renewal data AIC selects $K^*>0$ with asymptotic probability
\begin{equation}
P_{\rm AIC}=1-\exp\Big[-\sum_{k=1}^{K_{\max}}\frac1k\,P(\chi^2_k>2k)\Big],
\label{eq:woodroofe}
\end{equation}
equal to 0.248, 0.267, 0.284 and 0.288 for $K_{\max}=4$, 6, 12 and $\infty$, however long the data. This classical consequence of Spitzer's combinatorial lemma~\cite{Shibata1976,Woodroofe1982} makes AIC inconsistent for the order of a chain, and cross-validation with a fixed training fraction, asymptotically equivalent to it~\cite{Stone1977,Shao1993}, inherits the same bias. BIC is consistent~\cite{Katz1981,CsiszarShields2000}. Heavy tails play no role here, since in the log variable every sufficient statistic is bounded. We select the depth with BIC and calibrate every claim against surrogate data.

On walkers generated by the model itself with random parameters, the procedure does what it should (Appendix~\ref{app:selection}, \figref{fig:validation}). Planted and inferred couplings agree, and their differences, in units of the Fisher errors~\eqref{eq:fisher}, have mean $-0.01$ and standard deviation 0.97. BIC recovers the planted depth in 78\% of the walkers and never selects a deeper chain, whereas AIC does so in 24\%. On renewal walkers the false-positive rate of AIC stays at the level of Eq.~\eqref{eq:woodroofe} from 250 to 16\,000 flights, and that of BIC vanishes as the data grow.

\subsection{What else looks like memory}\label{sec:masquerade}

A real trajectory can mimic memory in ways that a planted walker cannot, and three of them matter here (Appendix~\ref{app:pitfalls}, \figref{fig:pitfalls}). Walkers differ: two flights of the same fast walker are both long, and a model with one field for all reads this as memory, at the largest depth allowed. One field per walker removes the effect, which is why we never pool walkers without their own fields. Flights are not observed but cut out of noisy tracks, and noise occasionally breaks a long flight into two long pieces, which manufactures a positive correlation; merging consecutive pieces that are nearly collinear undoes it, and a pipeline null, which rebuilds each track from its own flights in random order and analyzes it again, measures whatever memory the pipeline itself creates. Finally, the variable matters: the linear statistic $\ell_t\ell_{t-1}$ is dominated by a handful of extreme flights and misses a memory that the log statistic $x_tx_{t-1}$ detects. Different mechanisms of dependence also leave different signatures in how the selected depth grows with the length of the data, which helps to tell a finite memory from a slowly decaying one.

\subsection{Reading interactions from a group}\label{sec:groupavg}

For groups, three results, proved in Appendix~\ref{app:collective} and tested on synthetic collectives in Table~\ref{tab:collective}, organize the inference.

\emph{The mean-field measure is an exact likelihood.} The factorized measure of Ref.~\cite{Agliari2024}, evaluated on the data, equals the causal likelihood with $\theta_t=H+\sum_\tau J_\tau m_{t-\tau}$, as for single walkers in Sec.~\ref{sec:causal}. It can therefore be maximized exactly. The Gaussian autoregressive approximation used in the original implementation for convenience overestimates a planted coupling by 24\% in groups of five units; the exact estimator does not.

\emph{The group average contains own memory.} Let the units be independent, each with its own memory and autocovariance $C_i(\tau)$. A model of the group average then acquires delayed couplings, because the autocorrelation of the average contains the autocorrelation of each unit, Eq.~\eqref{eq:Rsplit}. Linearizing the inversion of Ref.~\cite{Agliari2024}, the fitted couplings converge to
\begin{equation}
\begin{gathered}
J^\star=\frac{\phi}{\Lambda'(\bar\theta)},\qquad\phi=\bar{\mathsf R}^{-1}\bar{\bm c},\\
\bar{\mathsf R}_{\tau\tau'}=\bar C(|\tau-\tau'|),\qquad \bar c_\tau=\bar C(\tau),
\end{gathered}
\label{eq:aggregation}
\end{equation}
the Yule--Walker coefficients of the mean individual autocovariance $\bar C=N^{-1}\sum_iC_i$, divided by the susceptibility of one unit. The factors of $N$ cancel, so these couplings do not shrink with the size of the group. In simulations without interactions the group-average analysis declares a delayed coupling in every data set, for $N=10$, 30 and 100 alike, with coefficients that follow Eq.~\eqref{eq:aggregation} lag by lag [\figref{fig:many}(b)]: $(0.468,0.225,0.177)$ against the predicted $(0.459,0.230,0.156)$ at $N=30$. These couplings are not an error of the analysis. They are the memories of the units, measured through their average, and a unit's own past is as legitimate a source of delayed correlation as the past of its neighbors; but the two are memories of a different nature, and the average cannot separate them. In the validation of Ref.~\cite{Agliari2024} on the Vicsek model, a peak at the first delay was correctly attributed to the inertia of each particle; Eq.~\eqref{eq:aggregation} shows that an inertia lasting longer than one frame produces peaks at longer delays as well. The self\,+\,social model assigns own memory to the couplings $a_\tau$: with regressors that leave out the unit itself, its social score vanishes exactly when units do not interact, and it declares a coupling in 1 of 120 interaction-free data sets.

\emph{Common drives look like interactions.} A field shared by all units, such as temperature, illumination or a pulse of chemokine, correlates units that never interact. Time fields $g_t$ absorb it, but they also absorb a global coupling: once each frame has its own field, the mean of the other units at an earlier frame is a function of that frame plus a multiple of the unit's own past, and global coupling and common drive become indistinguishable. A local coupling, which differs from unit to unit within a frame, survives. Local coupling with time fields is the only reading that declares no coupling under a drive alone and recovers a planted coupling without bias (Appendix~\ref{app:collective}, \figref{fig:collective}).

One last control is specific to real groups. In a time-shift surrogate each unit's track is displaced in time by an independent random offset. The surrogate keeps every unit's own memory, its heterogeneity and its marginal law, and it destroys every synchrony between units, both interactions and common drives. A coupling that survives in the surrogates is not an interaction.
